% SOURCE_STATEMENT: "Construct a self-adjoint Hamiltonian, bounded from below, that models pairwise point interactions in a many-body quantum system."
% Preparation and final self-audit: CONTRIBUTING.md prompts C and D; source/audit artifacts retained in the local physics_statements_2026_10_11 batch.
% !TeX program = lualatex
\documentclass[11pt]{article}
\usepackage[margin=25mm]{geometry}
\usepackage{fontspec}
\setmainfont{Latin Modern Roman}
\usepackage{amsmath,amssymb,mathtools}
\usepackage{unicode-math}
\setmathfont{Latin Modern Math}
\usepackage{xurl,hyperref}
\hypersetup{colorlinks=true,urlcolor=blue}
\setcounter{secnumdepth}{0}
\setlength{\parindent}{0pt}
\setlength{\parskip}{5pt}
\setlength{\emergencystretch}{3em}
\newcommand{\sref}[2]{\hyperlink{source:#1:#2}{[S#2]}}
\newcommand{\cataloguescope}[1]{\paragraph{Recorded target.}#1}
\newcommand{\R}{\mathbb R}
\newcommand{\C}{\mathbb C}
\newcommand{\N}{\mathbb N}
\newcommand{\Z}{\mathbb Z}
\newcommand{\Q}{\mathbb Q}
\newcommand{\F}{\mathbb F}
\newcommand{\eps}{\varepsilon}
\begin{document}
\section{Lower-Bounded Hamiltonians for Many-Body Point Interactions}\label{statement}

\subsection{Definitions and mathematical statement}

\paragraph{Shared notation.}$\mathbb N=\{1,2,\ldots\}$; the standard number fields have their usual meanings. All additional conventions are stated locally.


\paragraph{Definitions and conventions.}
A point interaction is a zero-range interaction specified by contact boundary conditions or an extension of a free differential operator, rather than a naive bounded multiplication potential. A Hamiltonian is self-adjoint on an explicitly stated dense domain; boundedness below means $\langle\psi,H\psi\rangle\geq-C\|\psi\|^2$ on its quadratic-form domain for some finite $C$. Pairwise contact conditions must be jointly compatible with the statistics, mass ratios and many-particle collision geometry.

\paragraph{Target retained from the canonical source catalogue.}
\paragraph{Catalogue wording retained for traceability.}
Construct a self-adjoint Hamiltonian, bounded from below, that models pairwise point interactions in a many-body quantum system.

\paragraph{Source scope and qualification.}
The report explains that unrestricted bosonic zero-range systems can suffer Thomas collapse and focuses on a fermion-plus-distinguishable-particle model. Thus this extraction is a construction question in admissible regimes, not an existence claim for every statistics or mass ratio.
\subsection{Short English statement}
Construct mathematically sound many-particle contact-interaction Hamiltonians that do not collapse to arbitrarily negative energy.
\subsection{Sources}
\begin{itemize}
\item[\hypertarget{source:statement:1}{S1}] ULAM.AI and the UnsolvedMath Contributors, UnsolvedMath, catalogue entry. Catalogue curation under CC BY 4.0; definitions and English exposition added locally.
\url{https://huggingface.co/datasets/ulamai/UnsolvedMath}.
\item[\hypertarget{source:statement:2}{S2}] Mathematical Questions and Challenges in Quantum Electrodynamics and its Applications (2017), Oberwolfach Reports, DOI 10.4171/owr/2017/41, pp. 2539--2599. Primary report accessed 2026-10-11; relevant PDF pages 40.
\url{https://doi.org/10.4171/owr/2017/41}.
\item[\hypertarget{source:statement:3}{S3}] Three-body zero-range interactions (catalogue-listed context; not independently consulted in this preparation).
\url{https://arxiv.org/abs/2107.07188}.
\item[\hypertarget{source:statement:4}{S4}] Contact-interaction Hamiltonians (catalogue-listed context; not independently consulted in this preparation).
\url{https://arxiv.org/abs/2403.12594}.
\end{itemize}
\label{end:statement}
\end{document}
